
\section{Harness Interface: Code for Reasoning, Acting, and
Environment Modeling}
\label{sec:foundations}

A harness turns a stateless language model into a functional 
agent by grounding its outputs in external execution, persistent 
state, and verifiable feedback. The most fundamental design 
question for any harness is therefore: \emph{what medium 
connects the model to its task environment?}

We argue that code is the answer. Unlike natural language, code 
is \emph{executable}, meaning model outputs become operations 
with formally verifiable outcomes; \emph{inspectable}, meaning 
intermediate computation is exposed as structured traces that 
the harness can read, store, and act upon; and \emph{stateful}, 
meaning the evolving program represents task progress in a 
persistent, modifiable form across steps. Crucially, these are 
not merely properties of code as a notation; they are properties 
that make code functional as a harness interface. Executability 
means the harness can verify what the model intended. 
Inspectability means failures can be diagnosed and fed back. 
Statefulness means the agent's interaction history is not lost 
between steps.


\paragraph{Scope boundary.}
We use \emph{code} broadly, but not metaphorically. In this
survey, code refers to executable or machine-checkable artifacts,
including programs, scripts, formal specifications, proof scripts,
API schemas, tool definitions, tests, repositories, simulators,
configuration files, and code-adjacent execution artifacts such as
traces and logs when they are produced by or consumed by executable
systems. By contrast, raw perception, physical state, human intent,
and model-internal latent reasoning are not themselves code.
They may be sensed, estimated, serialized, verified, or acted upon
through code, but they should not be conflated with the code
interface. This boundary is important because code as a harness
interface does not replace perception, embodiment, human goals, or
model reasoning; rather, it makes selected aspects of them
executable, inspectable, and stateful within the agent loop.

We organize this interface around three roles that code assumes in
agentic systems. \emph{Code for reasoning} externalizes internal
logic into verifiable computation, allowing external interpreters,
symbolic solvers, execution traces, or process rewards to check
and refine reasoning (\S\ref{subsec:reasoning}). \emph{Code for
acting} translates high-level intent into executable operations
grounded in embodied, GUI, software, or tool-use environments
(\S\ref{subsec:acting}). \emph{Code for environment modeling}
represents world state, transition dynamics, and feedback signals
through program states, repositories, simulators, tests, and logs
that agents can execute, edit, and query
(\S\ref{subsec:environment}). Overall, these roles define the
harness interface: code makes reasoning executable, action
programmable, and environment state inspectable.





\subsection{Code for Reasoning}
\label{subsec:reasoning}
% ------------------------------------------------------------

A central role of the agent harness is to transform model
reasoning from transient text generation into executable and verifiable computation. Early prompting techniques such as pure chain-of-thought (CoT)~\cite{wei2022chain} perform reasoning and computation
entirely in \textit{natural language}, forcing the model to both decompose problems and execute intermediate operations within a single latent textual process.
While language models are often effective at proposing reasoning steps, they remain unreliable at faithfully carrying out symbolic, logical, or arithmetic computation~\cite{gao2023pal}. More importantly, purely textual reasoning provides the agent harness
with little ability to verify intermediate states, inspect execution behavior, or persist computational progress across steps.

\textit{Code-for-reasoning} thus introduces code as the execution interface
between the model and the harness, moving beyond purely text-based reasoning. The model generates executable programs that external runtimes, interpreters, symbolic solvers, or verification modules can execute and evaluate.
This separates high-level reasoning from low-level computation:
the model proposes procedures, while the harness executes them,
observes runtime behavior, stores intermediate states, and feeds
execution results into future reasoning.
% As a result, reasoning becomes executable through program
% execution, inspectable through runtime traces, and stateful
% through persistent computational artifacts.
% Existing work can be organized into three paradigms:
% program-delegated reasoning, hybrid symbolic--neural execution,
% and iterative code-grounded reasoning. 

Recent work further broadens this interface from program execution as an external calculator to execution artifacts as reusable reasoning signals. Inputs and outputs, execution traces, variable states, control-flow structures, and function-level tests can all serve as intermediate states that the harness verifies, scores, and feeds back into subsequent reasoning. Existing work can therefore be organized into three paradigms: program-delegated reasoning, formal verification and symbolic reasoning, and iterative code-grounded reasoning.
We detail each of them in the following subsections.


\begin{figure}[t!]
    \centering
    \includegraphics[width=\linewidth]{figures/sec2teaser.pdf}
    \caption{Overview of code as the harness interface, connecting agents to reasoning, action, and environment modeling through executable programs, tool calls, state tracking, and feedback traces.}
    \label{fig:sec2}
\end{figure}


\begin{figure}[t!]
    \centering
    \includegraphics[width=0.85\linewidth]{figures/roadmap_sec2.pdf}
    \caption{Roadmap of the harness interface, organized by code's role in reasoning, acting, and environment modeling, with representative works ordered chronologically within each role.}
    \label{fig:roadmap_sec2}
\end{figure}


\subsubsection{Program-Delegated Reasoning}
Program-delegated reasoning uses executable programs as the
primary interface between problem decomposition and computation.
Instead of relying solely on natural language reasoning, the
model generates code that external interpreters execute to
produce formally grounded outputs.
Early works~\cite{nye2021show,gao2023pal} demonstrate that
delegating computation to programs substantially improves
reliability by moving intermediate reasoning into structured,
verifiable execution traces.
Program-of-Thoughts (PoT) prompting~\cite{chen2022program}
further systematizes this paradigm by explicitly decomposing
reasoning into executable programs, followed by extensions such
as POET~\cite{pi2022reasoning} and
MathCoder~\cite{wang2023mathcoder}, which improve execution
fidelity and domain specialization.
Subsequent work investigates the conditions under which program
delegation is effective, including the role of execution
correctness, task structure, and runtime interaction.
For example, Chain of Code (CoC)~\cite{li2023chain} and
CIRS~\cite{bi2024program} analyze how executable reasoning
changes failure modes relative to pure language-based reasoning.
Later directions extend this interface beyond isolated task
execution. Cross-lingual reasoning
frameworks~\cite{payoungkhamdee2025towards} demonstrate that
program-based reasoning can generalize across linguistic
environments through shared executable structure, while
method-based reasoning~\cite{su2025method} introduces reusable
programmatic procedures that persist across tasks.
More recent systems such as
CodeAdapt~\cite{zhang2025code} further suggest that tightly
coupling language models with executable reasoning interfaces can surpass specialized reasoning-oriented models. Additionally, CodeI/O~\cite{pmlr-v267-li25t} transforms contextually grounded programs into code input-output prediction tasks, exposing reasoning primitives such as logic-flow planning, state-space search, decision-tree traversal, and modular decomposition while preserving procedural rigor through executable verification.

\subsubsection{Formal Verification and Symbolic Reasoning Interfaces}

Hybrid neural-symbolic methods combine flexible language-based inference with structured symbolic computation, using code and symbolic artifacts as persistent intermediate representations rather than treating programs as mere generated text. Early formulations such as Graph-of-Thoughts~\cite{besta2024graph} generalize chain-of-thought reasoning into graph-structured trajectories, enabling intermediate states to branch, merge, and be reused. Building on this direction, self-verifying reflection~\cite{yu2025self}, MA-LoT~\cite{wang2025ma}, and Socratic self-refine~\cite{shi2025ssr} introduce iterative verification loops in which symbolic consistency checks guide the refinement of generated solution paths.

Recent work further tightens the coupling between neural generation and symbolic execution through code-based interfaces. CodeSteer~\cite{chen2025codesteer} and Code-as-Symbolic-Planner~\cite{chen2025code} explicitly coordinate free-form language reasoning with executable symbolic operations, treating programs as structured substrates that the harness can inspect, transform, and execute across multiple stages. VisualCoder~\cite{chi-etal-2025-visualcoder} extends this idea by making program behavior visible through control-flow representations. By aligning generated reasoning with visual control-flow graphs and execution paths, it turns dynamic program behavior into an inspectable artifact for program-behavior prediction. Together, these methods broaden the neural-symbolic interface from textual code to multimodal execution artifacts that a harness can reference, validate, and reuse.

A complementary line of work uses machine-verifiable formal languages as the reasoning interface itself. Proof assistants such as Lean~\cite{moura2021lean}, Isabelle~\cite{nipkow2002isabelle}, and Coq~\cite{barras1999coq} provide formal proof languages based on rigorous logical foundations, enabling each derivation step to be checked by a verifier. Early LLM-based theorem-proving systems, including ReProver~\cite{yang2023leandojo}, DeepSeek-Prover~\cite{xin2025deepseek}, and TheoremLlama~\cite{wang2024theoremllama}, establish practical recipes for combining language models with proof-assistant feedback in mathematical reasoning. More recent systems, such as DeepSeek-Prover-V2~\cite{ren2025deepseek2}, Kimina-Prover~\cite{wang2025kimina}, MA-LoT~\cite{wang2025ma}, and Goedel-Prover-V2~\cite{lin2025goedel2}, improve this process through deliberative proof search, self-correction, and repeated proof generation and verification.
Formal verification interfaces are also expanding beyond theorem proving in mathematics. HybridReasoning~\cite{wang2025let} applies formal provers to support natural-language reasoning; Lean4Physics~\cite{li2025lean4physics} and PhysLib~\cite{physlib} extend Lean-based verification to physics; and VERINA~\cite{ye2025verina} and Goedel-Code-Prover~\cite{li2026goedel} adapt formal methods to code verification. Lean4Agent~\cite{wang2026lean4agent} further extends this trajectory to agentic systems by using Lean4 to model and verify agent workflows and trajectories. From the harness perspective, these systems show how formal languages can serve not only as reasoning tools, but also as executable contracts that constrain, certify, and audit agent behavior.



\subsubsection{Iterative Code-Grounded Reasoning}
Iterative code-grounded reasoning focuses on closed-loop
interaction between generation, execution, and feedback.
In these systems, reasoning is not a single-pass process, but an
iterative computational trajectory grounded in executable state
transitions.
Early work such as NExT~\cite{ni2024next} trains models to
anticipate execution behavior by reasoning over program traces,
thereby grounding intermediate reasoning in runtime semantics.
Related efforts~\cite{armengol2025cannot} similarly emphasize
that executable traces provide a richer supervision signal than
final textual outputs alone.
Building on this foundation, subsequent approaches introduce
explicit generate--execute--verify--refine loops.
Methods such as CodePRM~\cite{li2025codeprm} and
ORPS~\cite{yu2024reasoning} use execution outcomes to evaluate
and refine intermediate reasoning trajectories, enabling the
harness to guide reasoning through runtime feedback rather than
pure next-token prediction.
Along the same direction, systems such as
CYCLE~\cite{ding2024cycle} and
Self-Edit~\cite{zhang2023self} iteratively revise generated
solutions using execution-aware correction signals.
Reinforcement learning further strengthens this paradigm by
treating execution feedback as an optimization signal over
reasoning trajectories.
Methods such as CodeRL~\cite{le2022coderl},
CodeRL+~\cite{jiang2025coderl+}, and
RLTF~\cite{liu2023rltf} optimize functional correctness through
unit-test-based rewards, while approaches such as
StepCoder~\cite{dou2024stepcoder} incorporate fine-grained
compiler and runtime feedback during optimization.
RLEF~\cite{gehring2024rlef} formalizes this interaction as
policy optimization grounded in multi-step execution feedback,
allowing reasoning policies to adapt through iterative runtime
interaction.
More recent approaches move toward fully interactive reasoning
environments.
For example, EG-CFG~\cite{lavon2025execution} injects execution
signals directly during generation to support step-level
correction, while systems such as
R1-Code-Interpreter~\cite{chen2025r1} interleave reasoning and
multiple rounds of code execution within persistent interactive
sessions.




\begin{table}[t]
\centering
\caption{
Representative systems where code serves as a reasoning substrate.
}
\renewcommand{\arraystretch}{1.15}
\setlength{\tabcolsep}{4pt}
\footnotesize
\begin{tabularx}{\textwidth}{p{2.8cm}p{2.2cm}p{3.1cm}X}
\toprule
\textbf{Method} & \textbf{Mechanism} & \textbf{Reasoning Paradigm} & \textbf{Key Innovation} \\
\midrule
PoT~\cite{chen2022program} & Delegated & Hybrid comments & Merges code with natural language CoT \\
PAL~\cite{gao2023pal} & Delegated & Program-aided & Decouples logic from computation \\
CodeAdapt~\cite{zhang2025code} & Delegated & Generalizable logic & Code-enabled LLMs outperforming reasoning models \\
CodeI/O~\cite{pmlr-v267-li25t} & Delegated & I/O prediction & Converts code into verifiable input-output reasoning tasks \\
SATLM~\cite{ye2023satlm} & Formal & SAT/SMT solving & Uses symbolic solvers as machine-checkable reasoning backends \\
ReProver~\cite{yang2023leandojo} & Formal & Lean proof search & Combines LLM generation with proof-assistant feedback \\
Dpsk-Prover~\cite{xin2025deepseek} & Formal & Lean theorem proving & Trains LLMs for formal mathematical proof generation \\
Dpsk-Prover-V2~\cite{ren2025deepseek2} & Formal & Deliberative proving & Lean proof search through decomposition and self-correction \\
Goedel-Code-Prover~\cite{li2026goedel} & Formal & Lean code proof & Searches hierarchical Lean proofs for code verification \\
Lean4Agent~\cite{wang2026lean4agent} & Formal & Agent verification & Models and verifies agent workflows and trajectories in Lean4 \\
Chain of Code~\cite{li2023chain} & Hybrid & LMulator & Simulates non-executable semantic code \\
SATLM~\cite{ye2023satlm} & Hybrid & Formal Logic & Uses SAT/SMT solvers as reasoning backend \\
CodeSteer~\cite{chen2025codesteer} & Hybrid & Symbolic control & Explicitly transitions between symbolic code and neural text \\
VisualCoder~\cite{chi-etal-2025-visualcoder} & Hybrid & CFG-grounded & Aligns code reasoning with visual control-flow artifacts. \\
NExT~\cite{ni2024next} & Iterative & Trace-grounded & Anticipates execution behavior via program traces \\
MathCoder~\cite{wang2023mathcoder} & Iterative & Feedback-driven SFT & Interleaves code, output, and reflection \\
CodePRM~\cite{li2025codeprm} & Iterative & Process rewards & Learns reward functions over reasoning-execution trajectories \\
RLEF~\cite{gehring2024rlef} & Iterative & Multi-step RL & Optimizes policy directly using execution feedback \\
EG-CFG~\cite{lavon2025execution} & Iterative & Execution-guided & Integrates execution signals directly during generation \\
R1-Code-Int.~\cite{chen2025r1} & Iterative & Fully interactive & Autonomously interleaves reasoning and multiple executions \\
ExecVerify~\cite{tang2026execverifywhiteboxrlverifiable} & Iterative & Stepwise RL & Uses statement- and variable-level execution rewards. \\
FunPRM~\cite{zhang2026funprmfunctionasstepprocessreward} & Iterative & Function-step PRM & Treats functions as verifiable process-reward units. \\
ReCode~\cite{fan2026recodereinforcingcodegeneration} & Iterative & Process RL & Reinforces code generation with reasoning-process rewards \\
\bottomrule
\end{tabularx}
\label{tab:code_reasoning_systems}
\end{table}


% ------------------------------------------------------------
\subsection{Code for Acting}
\label{subsec:acting}


Beyond reasoning, the agent must also connect the model to external environments where decisions produce real executable effects. At this stage, code no longer serves primarily as a medium for
computation, but as an action interface that converts model outputs into grounded operations such as tool invocations, robot-control policies, GUI actions, or software commands.
Through this interface, the harness translates high-level intent into executable behaviors that can interact with embodied, digital, and interactive environments.
The central challenge is therefore grounding: the harness must map abstract language outputs into executable behaviors that respect the constraints of the target environment, including embodiment limits, interface APIs, environment dynamics, and safety requirements. Unlike code-for-reasoning, where interpreters can often directly verify correctness, action execution occurs in partially observed and dynamically evolving environments, where failures may emerge through invalid state transitions, delayed feedback, or silent execution errors. For example, a robot may attempt to grasp an object outside its reachable workspace without producing an explicit runtime exception.


Importantly, executable action code is an interface to these
components, not a replacement for them. In embodied settings,
perception modules provide observations, affordance or feasibility
models estimate which actions are possible, motion planners and
controllers connect symbolic commands to sensors and actuators,
and safety layers constrain dangerous or invalid behavior. In GUI
and software settings, the analogous components include screen
parsers, DOM or accessibility trees, backend APIs, user-intent
models, permission systems, and programmatic validators. Code sits
between the model and these components: it serializes observations,
calls grounding and planning modules, invokes executable actions,
and exposes validation results back to the harness.

\textit{Code-for-acting} therefore introduces structured
executable programs as the control interface between the model
and the environment, allowing the harness to execute, monitor,
validate, reuse, and refine actions through interaction feedback.
This interface can be realized in different forms: a predefined
skill library, a generated control policy, a persistent skill
memory, a GUI/API tool protocol, or an explicit action-validation
harness. AutoHarness~\cite{lou2026autoharnessimprovingllmagents}
makes the last form explicit by automatically synthesizing a code
harness that mediates between the LLM and the environment,
filtering invalid actions before execution. This highlights the
core harness view of code-for-acting: code is not only the action
to be executed, but also the executable boundary that connects
model intent to perception, grounding, affordance estimates,
controllers, APIs, actuators, and safety constraints.




\begin{table}[t]
\centering
\caption{Representative systems where code serves as an action interface.}
\renewcommand{\arraystretch}{1.15}
\setlength{\tabcolsep}{4pt}
\footnotesize
\begin{tabularx}{\textwidth}{p{2.7cm}p{1.7cm}p{3.0cm}X}
\toprule
\textbf{Method} & \textbf{Mechanism} & \textbf{Action Paradigm} & \textbf{Key Innovation} \\
\midrule
AutoHarness~\cite{lou2026autoharnessimprovingllmagents} & Harness Gen. & Action validation & Synthesizes code harnesses that mediate model actions and filter invalid environment interactions \\
SayCan~\cite{ahn2022can} & Skill Selec. & Affordance-based & Links LLM plans to physical feasibility \\
KnowNo~\cite{ren2023robots} & Skill Selec. & Conformal prediction & Calibrates planner uncertainty for ambiguous instructions \\
SkillVLA~\cite{zhai2026skillvla} & Skill Selec. & Bimanual grounding & Extends grounding to combinatorial skill reuse \\
BOSS~\cite{zhang2023bootstrap} & Skill Selec. & Skill bootstrapping & Synthesizes new executable skill chains via guided practice \\
LLM-Guided Traj.~\cite{ha2023scaling} & Skill Selec. & Trajectory generation & Generates diverse manipulation trajectories and executable success conditions \\
LRLL~\cite{tziafas2024lifelong} & Skill Selec. & Lifelong grounding & Evolving skill interface via memory and self-exploration \\
CaP~\cite{liang2023code} & Policy Gen. & Hierarchical Python & Generates reactive robot control policies \\
RoboCodeX~\cite{mu2024robocodex} & Policy Gen. & Multimodal tree & Synthesizes tree-structured code across navigation \\
Code-BT~\cite{zhang2025codebt} & Policy Gen. & Behavior-tree & Imposes rule constraints via code-to-behavior-tree planning \\
ALRM~\cite{santos2026alrm} & Policy Gen. & Closed-loop control & Integrates programmatic generation with ReAct execution \\
CP-Agent~\cite{szeider2025cp} & Policy Gen. & Constraint solving & Uses persistent execution loops for formal constraint-model repair \\
Robot-Code Sim.~\cite{wang2025llm} & Policy Gen. & Static simulation & Uses LLMs as static simulators for robot code evaluation \\
GenSwarm~\cite{ji2026genswarm} & Policy Gen. & Multi-robot control & Coordinates policy generation and deployment across robotic agents \\
NormCode~\cite{guan2025normcode} & Policy Gen. & Governed interface & Enforces auditability and data isolation through semi-formal code \\
RACAS~\cite{ashley2026racas} & Policy Gen. & Cooperative control & Robot-agnostic architecture for closed-loop cooperative agents \\
Voyager~\cite{wang2023voyager} & Lifelong & Skill Library & Autonomous curriculum for open-ended tasks \\
LYRA~\cite{meng2025growing} & Lifelong & Human-in-loop & Encodes human corrections into reusable structured skills \\
ViReSkill~\cite{kagaya2025vireskill} & Lifelong & Vision-grounded & Replanning on failure using a skill-memory cache \\
UI-Voyager~\cite{lin2026ui} & Lifelong & Self-evolving & Rejection fine-tuning and self-distillation for mobile GUI agents \\
SkillsCrafter~\cite{wang2026lifelong} & Lifelong & Continual skills & Mitigates forgetting as executable manipulation skills accumulate \\
\bottomrule
\end{tabularx}
\label{tab:code_acting_systems}
\end{table}



\subsubsection{Grounded Skill Selection}
Grounded skill selection studies how the agent maps high-level language intent into executable behaviors through reusable skill interfaces. Rather than generating low-level actions directly, these systems treat the environment as a collection of executable capabilities
that the agent harness can invoke, compose, and refine under environmental constraints. SayCan~\cite{ahn2022can} establishes the core paradigm by coupling language planning with grounded skill execution, allowing the agent to select actions based not only on semantic relevance but also embodiment feasibility.
Subsequent work extends this execution interface in several directions. KnowNo~\cite{ren2023robots} introduces uncertainty-aware control through conformal prediction, enabling the harness to detect ambiguous states and trigger clarification before unsafe
execution. BOSS~\cite{zhang2023bootstrap} addresses the rigidity of fixed skill libraries by using language-guided practice to synthesize
new executable skill chains, allowing the harness to expand its
action space over time.
Similarly, \cite{ha2023scaling} tackles the data bottleneck of
grounded interaction by using LLM-guided generation to construct
diverse manipulation trajectories and executable success
conditions for automatic retry and relabeling.
Beyond static execution, LRLL~\cite{tziafas2024lifelong}
introduces memory and self-guided exploration to maintain a
persistent and evolving skill interface across tasks.
Finally, SkillVLA~\cite{zhai2026skillvla} extends this paradigm
to combinatorial bimanual interaction, emphasizing that grounded
action interfaces must support structured skill reuse and
recomposition under increasingly complex embodiment settings.

\subsubsection{Programmatic Policy Generation}
Programmatic policy generation treats code itself as the control
interface between the model and the environment.
Instead of selecting from predefined skills, the harness directly
materializes executable policies as programs that specify control
logic, perception-conditioned branching, feedback loops, and API
interaction.
CaP~\cite{liang2023code} crystallizes this paradigm by framing
LLM-generated Python programs as executable robot policies.
Building on this idea, RoboCodeX~\cite{mu2024robocodex}
introduces multimodal and tree-structured code generation to
support more complex manipulation and navigation behaviors.
Subsequent work focuses on scaling the interaction substrate.
RoboPro~\cite{xie2025robotic} synthesizes executable policy code
from large-scale in-the-wild videos, while
Code-BT~\cite{zhang2025codebt} compiles generated programs into
behavior-tree controllers that support constrained execution and
iterative runtime feedback.
Beyond robotics, CP-Agent~\cite{szeider2025cp} demonstrates that
persistent execution loops can support formal constraint-solving
agents through iterative execution and repair.
To reduce dependence on expensive physical environments,
\cite{wang2025llm} configures language models as static execution
simulators for robot code evaluation.
GenSwarm~\cite{ji2026genswarm} further extends programmatic
control to multi-agent robotic systems, where the harness must
coordinate policy generation, constraint analysis, and deployment
across multiple embodied agents.
At the systems level, NormCode~\cite{guan2025normcode}
emphasizes governance and auditability by introducing a
semi-formal programming interface with enforced data isolation,
allowing execution traces and control logic to remain
inspectable and constrained.
Finally, ALRM~\cite{santos2026alrm} and
RACAS~\cite{ashley2026racas} consolidate these ideas into
persistent closed-loop control architectures that integrate code
generation, execution, monitoring, and iterative interaction
within unified agent harnesses.

\subsubsection{Lifelong Code-Based Agents}
Lifelong code-based agents study how executable interaction
interfaces can persist, evolve, and accumulate capabilities over
long-horizon interaction.
In these systems, code is not only an execution mechanism, but
also a persistent memory substrate through which the harness
stores reusable behaviors, interaction traces, and environment
knowledge.
Voyager~\cite{wang2023voyager} establishes this paradigm through
an automatic curriculum and continually expanding executable
skill library for open-ended interaction in Minecraft.
Extending this idea to embodied environments,
LRLL~\cite{tziafas2024lifelong} introduces persistent memory,
self-guided task exploration, and skill abstraction to overcome
the limitations of fixed policy libraries without requiring
gradient updates.
A central challenge in lifelong harnesses is that interaction
feedback and corrections are often transient and difficult to
reuse.
LYRA~\cite{meng2025growing} addresses this issue by converting
human corrections into reusable executable skills and
retrieval-augmented memory structures.
Similarly, ViReSkill~\cite{kagaya2025vireskill} combines
vision-grounded replanning with skill-memory caching to maintain
stable interaction under environmental failures and output
variability.
Recent work further focuses on continual adaptation and
self-evolution under persistent deployment.
SkillsCrafter~\cite{wang2026lifelong} introduces continual
language-conditioned manipulation structures to mitigate
catastrophic forgetting as executable capabilities accumulate,
while UI-Voyager~\cite{lin2026ui} generalizes the self-evolving
interaction paradigm to GUI agents through failure-driven
adaptation and self-distillation.
Together, these systems move beyond one-shot execution toward
persistent agent harnesses that continuously expand, refine, and
reuse executable interaction interfaces over time.

% ------------------------------------------------------------
\subsection{Code for Environment}
\label{subsec:environment}
% ------------------------------------------------------------


The agent must also maintain an explicit representation
of the environment with which the agent interacts.
Without such a representation, the environment is exposed to the agent only indirectly through textual observations, API returns, or sparse feedback signals.
As a result, environment state often remains implicit, transient, and difficult to verify, making it challenging to
track state transitions, evaluate interaction outcomes, or reuse
past interaction history across long-horizon tasks.
This limitation becomes particularly severe in complex software,
robotic, and multi-step interactive environments, where
successful interaction depends on maintaining consistent world
state and grounded feedback over time.


\textit{Code-for-environment} addresses this limitation by
introducing executable programs as the environment interface
itself.
Instead of treating the environment as an opaque external process,
these systems materialize environment structure and dynamics
through computational artifacts such as simulators,
repositories, tests, execution traces, logs, and state-transition
programs.
This allows the agent to explicitly store, inspect, execute,
and modify environment state throughout interaction.
% Consequently, environment dynamics become executable through
% runtime transitions, inspectable through structured traces, and
% stateful through persistent computational representations. 
Representing environments through executable code provides two
major advantages.
First, executable environments expose verifiable state
transitions, allowing the agent to evaluate interaction
outcomes through execution rather than ambiguous natural-language
judgment.
Second, code-based environments are persistent and
modifiable that agents can query, simulate, edit,
and refine during interaction.
Rather than interacting with an opaque world solely through language, agent harness can ground reasoning and action in explicit computational state and runtime dynamics. 
% Existing work in this direction can be organized into three paradigms: structured world representations, execution-trace world modeling, and code-grounded evaluation environments. 
Existing work in this direction can be organized into four
paradigms: structured world representations, execution-trace world
modeling, code-grounded evaluation environments, and verifiable
environment construction.
% We introduce these paradigms in the following subsections.

\begin{table}[t]
\centering
\caption{Representative systems where code serves as an environment representation.}
\renewcommand{\arraystretch}{1.15}
\setlength{\tabcolsep}{4pt}
\footnotesize
\begin{tabularx}{\textwidth}{p{2.8cm}p{2.0cm}p{3.1cm}X}
\toprule
\textbf{Method} & \textbf{Mechanism} & \textbf{Environment Paradigm} & \textbf{Key Innovation} \\
\midrule
ViStruct~\cite{chen2023vistruct} & Structured & Class/object hierarchy & Encodes visual scenes as data structures \\
FactoredScenes~\cite{hsu2025programs} & Structured & Room programs & Composes object/relation functions for 3D layout generation \\
PoE-World~\cite{piriyakulkij2025poe} & Structured & Programmatic experts & Scales symbolic world models beyond simple grid-worlds \\
Code2World~\cite{zheng2026code2world} & Structured & Render-aware RL & Re-frames GUI state prediction as renderable HTML generation \\
SemCoder~\cite{ding2024semcoder} & Trace-based & Semantic alignment & Pairs code with detailed execution traces \\
WorldCoder~\cite{tang2024worldcoder} & Trace-based & Model-based RL & Synthesizes transition and reward models \\
CWM~\cite{copet2025cwm} & Trace-based & Open-weights trace & Trains large LLMs natively on program execution traces \\
RWML~\cite{yu2026reinforcement} & Trace-based & Self-supervised RL & Aligns simulated next states with realized environment states \\
AWM~\cite{wang2026agent} & Trace-based & World-modeling & Aligns multiple executable world models across tasks \\
WorldMind~\cite{ren2026aligning} & Trace-based & Model fusion & Coordinates executable world models from knowledge sources \\
SWE-bench~\cite{jimenez2023swe} & Evaluation & Repo-level testing & Uses unit tests as objective world states \\
AgentBench~\cite{liu2023agentbench} & Evaluation & Multi-env interaction & Benchmarks across OS, databases, and games \\
CRUXEval~\cite{gu2024cruxeval} & Evaluation & Execution tasks & Benchmarks functional input and output prediction \\
End Terms.~\cite{gandhi2026endless} & Evaluation & Procedural RL envs & Automates generation of terminal-use evaluation tasks \\
InterCode~\cite{yang2023intercode} & Evaluation & Interactive execution & Frames coding tasks as actions with sandbox feedback \\
LiveCodeBench~\cite{jain2024livecodebench} & Evaluation & Live coding eval & Continuously updates execution-based evaluation pipelines \\
CRUXEval-X~\cite{xu2025cruxeval} & Evaluation & Multilingual execution & Extends input-output execution evaluation across languages \\
CoRe~\cite{xie2025core} & Evaluation & Runtime reasoning & Evaluates code reasoning through execution-centered tasks \\
CodeGlance~\cite{wang2026codeglance} & Evaluation & Multimodal code eval & Evaluates code understanding under visual and structural settings \\
SWE-smith~\cite{yang2025swesmithscalingdatasoftware} & Construction & Synthetic SWE envs & Generates repository-level tasks and execution environments \\
EnvScaler~\cite{song2026envscalerscalingtoolinteractiveenvironments} & Construction & Tool-interactive envs & Synthesizes tool-use environments with programmatic validators \\
\bottomrule
\end{tabularx}
\label{tab:code_environment_systems}
\end{table}



\subsubsection{Structured World Representations}
Structured world representations model environments through
explicit programmatic structures that the agent can execute,
inspect, and manipulate.
Rather than representing the environment solely through latent
embeddings or textual descriptions, these approaches encode world
state, object relations, spatial layouts, and interaction
dynamics as structured computational artifacts.
For example, ViStruct~\cite{chen2023vistruct} uses
programming-language structure as an explicit interface for
visual structural knowledge extraction, enabling multi-granular
visual events to be represented through consistent executable
structures.
FactoredScenes~\cite{hsu2025programs} similarly models indoor
environments as compositional ``room programs,'' where reusable
object and relation functions define physically consistent scene
layouts.
Extending this idea to scalable symbolic world modeling,
PoE-World~\cite{piriyakulkij2025poe} introduces a compositional
framework that combines many small programmatic experts to
represent increasingly complex environment dynamics.
More recent systems broaden structured environment interfaces to
high-fidelity interactive worlds.
Code2World~\cite{zheng2026code2world} reframes GUI state
prediction as renderable HTML generation, allowing environment
transitions to be represented and evaluated through executable
rendering code.
Code2Worlds~\cite{zhang2026code2worlds} further extends this
paradigm to 4D simulated environments through language-to-simulation
program generation, where physics-aware execution loops
reduce semantic-physical inconsistencies during environment
construction and interaction.

\subsubsection{Execution-Trace World Modeling}
Execution-trace world modeling studies how the agent can learn
environment dynamics directly from executable interaction traces.
Instead of treating execution merely as a final evaluation step,
these approaches model runtime transitions themselves as the
primary representation of environment behavior.
SemCoder~\cite{ding2024semcoder} bridges static programs and
runtime semantics by training language models to reason about
functional behavior, statement-level execution effects, and
input-output transitions.
Building on this perspective, Code World
Model~(CWM)~\cite{copet2025cwm} learns predictive world models
directly from program traces, enabling the agent to anticipate
future environment states through executable dynamics.
WorldCoder~\cite{tang2024worldcoder} further introduces a
model-based interaction framework in which the agent explicitly
writes and updates executable world models represented as Python
programs.
Rather than storing environment knowledge implicitly in model
parameters alone, the agent maintains editable computational
representations that can be executed, revised, and reused during
planning and interaction.
Subsequent work extends this paradigm toward continual and
interactive world-model adaptation.
RWML~\cite{yu2026reinforcement} combines execution traces with
reinforcement learning to refine environment dynamics through
runtime interaction, while
AWM~\cite{wang2026agent} and
WorldMind~\cite{ren2026aligning} study how multiple executable
world models can be aligned, fused, and coordinated across tasks
and knowledge sources.

\subsubsection{Code-Grounded Evaluation Environments}
Code-grounded evaluation environments use executable systems as
the interface for measuring agent behavior and interaction
quality.
Unlike static benchmarks based solely on textual outputs, these
environments expose explicit runtime state transitions, execution
feedback, and verifiable interaction outcomes that the agent
can directly observe and evaluate.
InterCode~\cite{yang2023intercode} establishes this paradigm by
reframing coding tasks as interactive execution environments,
where code acts as actions, execution feedback serves as
observations, and sandboxed runtimes provide grounded
interaction.
CRUXEval~\cite{gu2024cruxeval} further evaluates program
understanding through executable input-output prediction tasks,
while LiveCodeBench~\cite{jain2024livecodebench} introduces
continuously updated evaluation pipelines that assess execution,
self-repair, and runtime reasoning capabilities under evolving
problem distributions.
SWE-bench~\cite{jimenez2023swe} extends executable evaluation to
real-world software repositories, where agents must modify
large-scale codebases and are evaluated through repository-level
unit-test execution rather than textual correctness alone.
More broadly, AgentBench~\cite{liu2023agentbench} demonstrates
that executable interaction environments can evaluate reasoning
and decision-making across diverse embodied and digital tasks.
Subsequent benchmarks such as
CRUXEval-X~\cite{xu2025cruxeval},
CoRe~\cite{xie2025core},
GeoGramBench~\cite{luo2025geogrambench},
CodeGlance~\cite{wang2026codeglance}, and
Endless Terminals~\cite{gandhi2026endless} further expand this
paradigm toward multilingual, multimodal, and continuously
interactive evaluation settings, where runtime interaction rather
than static answer matching becomes the primary evaluation
interface.



\subsubsection{Verifiable Environment Construction}
A newer direction treats executable environments not only as
benchmarks to evaluate agents, but as harness artifacts that can
be synthesized, scaled, and validated programmatically. This is
especially important for long-horizon agents, where the harness
must provide not only a task prompt, but also a runnable state,
transition dynamics, feedback channels, and verification oracles.
SWE-smith~\cite{yang2025swesmithscalingdatasoftware} scales
software-engineering agent data by constructing repository-level
tasks and execution environments from existing codebases, turning
software repositories into reproducible program worlds for agent
training and evaluation. EnvScaler~\cite{song2026envscalerscalingtoolinteractiveenvironments}
extends this idea beyond software engineering by programmatically
synthesizing tool-interactive environments together with scenarios
and rule-based trajectory validators. From the harness perspective,
these methods make the environment interface itself an object of
construction: code specifies not only what the agent edits or
executes, but also the state transitions, tool affordances, and
verifiers that determine whether an interaction has succeeded.


